\section{Online, reference-free signals}
\label{sec:method}

Every signal below is a function of the window $[t-w+1,t]$ and of statistics accumulated on steps
$<t-w+1$. All are computed by the released implementation (\code{src/features.py}) from normalized
actions and extracted entities, so the same feature vector can be produced by any runtime that can log
tool calls and their outputs.

\subsection{Action normalization}
Raw tool calls differ across scaffolds (a shell command, a \code{str\_replace} call, a \code{Grep} tool),
so we first canonicalize each call into a tuple: a tool class (\code{shell}, \code{read}, \code{edit},
\code{search}, \code{verify}, \code{install}, \code{web}, \code{control}), a leading verb, a set of
plausible file or symbol targets, search terms, and a coarse intent (\code{explore}, \code{read},
\code{edit}, \code{verify}, \code{install}, \code{control}). Compound shell commands are split on pipes and
redirections and classified by priority, so \code{cd /app \&\& pytest -x} is a verification, not a
directory change, and multi-tool steps are flattened into an action sequence. This normalization is what
makes ``the same action again'' comparable across scaffolds; without it, repetition detection is
scaffold-specific.

\subsection{Evidence: entities and verification state}
From each observation we extract typed entities: file paths (including traceback frames and
\code{[File:...]} headers), function and symbol names, exception classes, normalized exception messages,
test identifiers, and metric-like key/value pairs, de-duplicated per step and identified by
\code{kind:value} with paths normalized to repository-relative form. In parallel we read \emph{objective
state} out of each observation: process exit codes, test pass/fail/error counts, build success or
failure markers, crash and timeout markers, and the terminal error signature (exception class plus
normalized message). Comparing two consecutive verification states yields the verification delta.

\subsection{Task grounding without a model}
The only task-specific input is the task statement. We tokenize it into content words and score an
entity's relevance as the weighted fraction of its subtokens appearing in the statement, with weights
being inverse-document-frequencies estimated on the training split (so a token like \code{file}
contributes little and \code{pyproject} contributes substantially). This is a zero-cost relevance prior:
no second model, no reference solution, no hand-written rules per task. It is deliberately crude, and
Section~\ref{sec:results} shows both that it is enough to change the operating point and where it fails.

\subsection{Six channels}
\begin{description}[leftmargin=1.1em,style=nextline,itemsep=2pt]
\item[\rep{} --- repetition.] Exact-duplicate fraction of normalized action signatures, at target level
and at command-family level; the longest exact repetition; the fraction of window actions already seen
anywhere earlier; and a periodicity statistic for $2$-, $3$- and $4$-cycles.
\item[\nov{} --- relevance-blind novelty.] Rate of entities never observed before, split by kind, plus
the fraction of distinct action signatures and the mean observation length. This is the ``is anything
new happening?'' channel, blind to whether the novelty matters.
\item[\evid{} --- task-grounded evidence.] Novel entities weighted by relevance: rate of new relevant
entities, share of new entities that are relevant, rate of high-relevance novelty, relevance-weighted
novelty mass, and a persistence term for novel entities that reappear inside the window (transient noise
does not count).
\item[\ver{} --- verification deltas.] Count of verification events, improvements and regressions versus
the previous verification, net progress, steps since the last improvement, rate of \emph{new} error
signatures, test-count deltas, and the modal error signature's dominance.
\item[\work{} --- workspace dynamics.] Number of edits, share of edits aimed at previously untouched
targets, churn (repeat edits to the same target), re-edit rate, and edits performed after the agent
declared completion.
\item[\sem{} --- semantic redundancy.] Rolling diversity of step representations: $1-$ mean pairwise
cosine similarity inside the window, similarity to the nearest earlier step, distance from the window
centroid, and novelty against the whole prefix. This is the livelock-style channel;
Section~\ref{sec:results} documents that in this environment the representation had to be a hashing
bag-of-words surrogate rather than a pretrained encoder, because the model host was unreachable.
\end{description}

\subsection{Monitors}
A monitor maps a feature vector to a stagnation score in $[0,1]$ and raises an alarm when the score stays
above a threshold for two consecutive steps, subject to a minimum step margin. We evaluate three families:
fixed baselines (B1 alarms after a step budget of 30 or 60 steps; B2 alarms on $k$ identical normalized
actions in a row, $k=3,5$); hand-designed channel monitors (signed, standardised sums of one channel's
features or of a chosen combination, squashed through a logistic, with standardisation statistics from the
training split); and learned monitors ($\ell_2$-regularised logistic regression on a channel's features,
fitted on annotated training windows only with balanced class weights). We train one learned monitor per
channel and one per combination --- evidence, evidence$+$verification, evidence$+$novelty,
evidence$+$semantic, all channels --- so the ablation is learned rather than hand-tuned.

\subsection{Operational alarm metrics}
Window-level AUC tells us how well a monitor ranks, but a runtime decides \emph{when} to act, so we also run
each monitor sequentially over each test trajectory and evaluate the first alarm. \textbf{Detection}: the
alarm lands inside an annotated stagnant region. \textbf{False stop}: the alarm lands inside an annotated
productive region, or before the first stagnant region while productive work remained; because a false stop
destroys work the annotators called productive, we count it as \emph{negative} savings. \textbf{Savings}:
the signed number of steps between the alarm and the end of the run, whose oracle bound is the step index
of the first annotated stagnant region. \textbf{Latency}: steps between the start of the stagnant region
and the alarm. All headline numbers are reported at fixed false-stop budgets (1\%, 2\%, 5\%, 10\%, 20\%),
following the practice of early-termination work, and the monotone curve of savings against false-stop rate
is reported in full rather than at a single cherry-picked threshold.
